\documentclass[11pt]{article}
\usepackage[T1]{fontenc}
\usepackage{lmodern}
\usepackage[a4paper,margin=28mm]{geometry}
\usepackage{amsmath,amssymb,amsthm,mathtools}
\usepackage[hidelinks]{hyperref}
\newtheorem{theorem}{Theorem}[section]
\newtheorem{lemma}[theorem]{Lemma}
\newtheorem{proposition}[theorem]{Proposition}
\theoremstyle{remark}
\newtheorem{remark}[theorem]{Remark}
\newcommand{\dd}{\,d\mu}
\newcommand{\E}{\mathcal E}
\newcommand{\norm}[1]{\left\lVert#1\right\rVert}
\DeclareMathOperator{\Ric}{Ric}
\DeclareMathOperator{\Vol}{Vol}
\DeclareMathOperator{\supp}{supp}
\title{Point-source solutions for fractional porous medium diffusion on complete manifolds}
\author{Kaiwen Jin\\\small\texttt{yui.ui.siki@gmail.com}}
\date{}
\begin{document}
\maketitle
\begin{abstract}
We construct solutions of the spectral fractional porous medium equation
starting from a point mass on a complete noncompact Riemannian manifold.
The assumptions are a lower Ricci curvature bound and a uniform
Euclidean Faber--Krahn inequality. Every exponent $m>1$, fractional order
$0<s<1$, prescribed source point and positive mass is allowed. The solution
is strongly continuous in $L^1$ at positive times, locally bounded, and
has the prescribed initial trace against every bounded continuous function.
The construction combines the known smoothing estimates for integrable
data with averages of the linear heat semigroup. These averages have small
fractional Laplacian uniformly in their spatial cutoff, which controls mass
at infinity while passing to a point-source limit.
\end{abstract}

\section{Introduction and main result}

The porous medium equation with fractional diffusion describes slow
nonlocal spreading of a concentration. A point-source solution starts with
all the material concentrated at one location. On a manifold the geometry
affects the diffusion, and the existence question can be asked without a
radial coordinate or a spatial scaling symmetry.

For the Euclidean fractional porous medium equation, finite nonnegative
measure data are covered by Grillo, Muratori and Punzo
\cite[Theorem~3.2]{GMP}. Berchio, Bonforte, Grillo and Muratori
\cite{BBGM} developed a theory on complete noncompact manifolds with a
Ricci lower bound and a uniform Faber--Krahn inequality. Their construction
gives weak dual solutions for integrable functions and for a larger
weighted class, together with smoothing estimates. Existence for a Dirac
initial datum is listed among the questions in \cite[Section~7]{BBGM}
and in Grillo's 2025 workshop report \cite{OWR}. Here we prove existence
in the distributional class stated below.

Throughout, $(M,g)$ is a smooth complete, connected, noncompact Riemannian
manifold of dimension $N\ge2$, with Riemannian volume $\mu$. Assume that
\begin{equation}\label{eq:geometry}
 \Ric\ge-(N-1)k\quad\text{for some }k\ge0,
 \qquad \lambda_1(D)\ge c\Vol(D)^{-2/N}
\end{equation}
for some $c>0$ and every relatively compact smooth domain $D\subset M$.
Here $\lambda_1(D)$ is the first Dirichlet eigenvalue of $-\Delta_M$.
Let $L$ be the nonnegative Friedrichs realization of $-\Delta_M$ on
$L^2(M)$ and let $B=L^s$, defined by the spectral calculus, where $0<s<1$.
We consider
\begin{equation}\label{eq:pde}
 \partial_tu+B(u^m)=0,\qquad m>1.
\end{equation}
The spaces $L^p(M)$ use $\mu$, and $C_b(M)$ denotes bounded continuous
functions. The distributional equation means
\begin{equation}\label{eq:weak}
 \int_0^\infty\!\int_M u\,\partial_t\phi\dd\,dt
 =\int_0^\infty\!\int_M u^m B\phi\dd\,dt,
 \qquad \phi\in C_c^\infty((0,\infty)\times M),
\end{equation}
with an absolutely integrable right side.

\begin{theorem}\label{thm:main}
Assume \eqref{eq:geometry}. For every $m>1$, $0<s<1$, $o\in M$ and
$A>0$, there is a nonnegative solution of \eqref{eq:weak} such that
\begin{equation}\label{eq:regularity}
 u\in C((0,\infty);L^1(M))\cap L^\infty_{\mathrm{loc}}((0,\infty)\times M),
 \qquad u^m\in L^1_{\mathrm{loc}}((0,\infty)\times M),
\end{equation}
and, for every $T>0$,
\begin{equation}\label{eq:mass}
 \sup_{0<t<T}\norm{u(t)}_1=A>0.
\end{equation}
Its initial trace is $A\delta_o$ in the sense that
\begin{equation}\label{eq:trace}
 \lim_{t\downarrow0}\int_Mu(t)\psi\dd=A\psi(o)
 \qquad\text{for every }\psi\in C_b(M).
\end{equation}
\end{theorem}

The assertion concerns existence on each manifold satisfying
\eqref{eq:geometry}. Uniqueness and a common large-time profile are separate
questions. The proof uses the integrable-data semigroup and smoothing
estimate of \cite{BBGM}. The additional step is to pass from concentrating
initial functions to a curve in the global space $L^1(M)$.

To do this, we first obtain local compactness at positive times from an
energy estimate for $u^m$. Averages of the linear heat semigroup then give
tightness over each finite time interval. At suitable positive times these
two facts yield convergence in $L^1(M)$, and the contraction of the same
nonlinear semigroup propagates that convergence to all later times. Finally,
the smoothing estimate makes the nonlinear flux integrable down to time
zero, which identifies the trace. The argument does not use a Green-potential
uniqueness theorem.

\section{Bounded-data solutions and their estimates}

Write $P_r=e^{-rL}$ for the linear heat semigroup. Under
\eqref{eq:geometry}, its symmetric kernel $p_r(x,y)$ satisfies
\begin{equation}\label{eq:heatkernel}
 0\le p_r(x,y)\le C r^{-N/2}
       \exp\!\left(-\frac{d(x,y)^2}{Cr}\right),\qquad r>0,
\end{equation}
by \cite[Assumption~2.1 and (2.3)]{BBGM}. Constants denoted by $C$ may
change from line to line. The Ricci bound also implies stochastic
completeness, so $P_r1=1$; see \cite[Theorem~15.4(a)]{Gri}. For example,
the radial curvature condition there holds with
$q(r)=\sqrt{k}+1+r$, and $\int^\infty dr/q(r)=\infty$.

The Faber--Krahn assumption in \cite{BBGM} is stated for arbitrary
relatively compact open sets. This follows from \eqref{eq:geometry} by
the Rayleigh characterization. Indeed, enclose the support of a smooth
compactly supported test function in finitely many smooth components
$D_i$ inside the given open set $\Omega$. Applying
\eqref{eq:geometry} on each component gives
\[
 \int_\Omega |\nabla\varphi|^2\dd
 \ge \sum_i c\Vol(D_i)^{-2/N}\int_{D_i}|\varphi|^2\dd
 \ge c\Vol(\Omega)^{-2/N}\int_\Omega|\varphi|^2\dd.
\]
Thus the source assumptions apply also with the smooth-domain convention
used here. When $k=0$, a larger positive Ricci-bound constant meets the
source's convention.

We use the following results for bounded integrable data. Set
\begin{equation}\label{eq:exponents}
 \alpha=\frac{N}{2s+N(m-1)},\qquad
 \beta=\frac{2s}{2s+N(m-1)}.
\end{equation}

\begin{lemma}[The bounded-data semigroup]\label{lem:semigroup}
For every nonnegative $v\in L^1(M)\cap L^\infty(M)$ there is a mild
solution $S_tv\in C([0,\infty);L^1(M))$ of \eqref{eq:pde}.
These solutions form a positivity-preserving semigroup and, for every
pair of such data $v,z$, satisfy
\begin{equation}\label{eq:contraction}
 \norm{S_tv-S_tz}_1\le\norm{v-z}_1,
 \qquad \norm{S_tv}_p\le\norm v_p\quad(1\le p\le\infty),
\end{equation}
and obey
\begin{equation}\label{eq:smoothing}
 \norm{S_tv}_\infty\le C_0
   \max\{t^{-\alpha}\norm v_1^\beta,\norm v_1\},\qquad t>0,
\end{equation}
where $C_0=C_0(N,k,c,s,m)$. Their implicit Euler approximations satisfy
\begin{equation}\label{eq:euler}
 z_j+hB(z_j^m)=z_{j-1},\qquad z_j^m\in D(B\text{ on }L^2(M)),
\end{equation}
and converge uniformly on each finite time interval in $L^1(M)$.
\end{lemma}

\begin{proof}
The semigroup, contraction, norm estimates and Euler construction are
\cite[Proposition~5.1 and (5.5)--(5.7)]{BBGM}. In particular, the proof
there establishes the stated $L^2$ operator domain after a positive
zeroth-order regularization. Proposition~5.3 identifies these mild
solutions as weak dual solutions. The construction in Section~6.1 starts
from $\mathbf1_{B_j}(v\wedge j)$, where $B_j$ is the geodesic ball of
radius $j$ about a fixed point. For bounded compactly supported $v$
this sequence is eventually $v$. For general bounded integrable $v$,
the contraction identifies the limit with $S_tv$.
Consequently \cite[Theorem~2.6, (2.13)]{BBGM} gives
\eqref{eq:smoothing} for this same semigroup. The numbered source
statements here refer to the version specified in the bibliography.
\end{proof}

For the energy and time estimates, put
\[
 V=D(L^{s/2}),\qquad
 \norm v_V^2=\norm v_2^2+\norm{L^{s/2}v}_2^2,
 \qquad \E(v,z)=\langle L^{s/2}v,L^{s/2}z\rangle_2.
\]
The form defines $B:V\to V^*$ using the $L^2$ pivot space. On its
operator domain this is the spectral operator already defined above.

\begin{lemma}[Energy and scalar testing]\label{lem:energy}
Let $u(t)=S_tv$ for nonnegative $v\in L^1\cap L^\infty$. For
$0\le a<b<\infty$,
\begin{equation}\label{eq:energy}
 \int_a^b\E(u^m,u^m)\,dt+
 \frac{\norm{u(b)}_{m+1}^{m+1}}{m+1}
 \le\frac{\norm{u(a)}_{m+1}^{m+1}}{m+1}.
\end{equation}
In particular $u^m\in L^2((a,b);V)$ and
$\partial_tu=-B(u^m)$ in $L^2((a,b);V^*)$.
For every $f\in D(B)\cap L^\infty(M)$ such that $Bf\in L^\infty(M)$,
\begin{equation}\label{eq:scalar}
 \langle u(t)-v,f\rangle
   =-\int_0^t\langle u(r)^m,Bf\rangle\,dr,\qquad t\ge0.
\end{equation}
\end{lemma}

\begin{proof}
The domain assertion in \eqref{eq:euler} permits an $L^2$ pairing with
$z_j^m$. Convexity gives
\[
 \frac{\norm{z_j}_{m+1}^{m+1}-\norm{z_{j-1}}_{m+1}^{m+1}}{m+1}
       +h\E(z_j^m,z_j^m)\le0.
\]
For this fixed datum the Euler curves have a common bound
$K_v=\norm v_\infty$. Uniform $L^1$ convergence and
$\int|z-u|^p\le(2K_v)^{p-1}\norm{z-u}_1$ give convergence in every finite
$L^p$, uniformly on bounded time intervals. The powers therefore converge
strongly in space-time $L^2$. Their form-energy bound gives a weak limit
in $L^2((a,b);V)$, necessarily $u^m$. Lower semicontinuity proves
\eqref{eq:energy}, first from time zero and then after a semigroup restart.
The $L^2$ part of the $V$ norm is bounded by
$\norm{u(t)^m}_2^2\le\norm{u(t)}_1\norm{u(t)}_\infty^{2m-1}$.

The piecewise affine Euler interpolation has derivative $-B(z_j^m)$.
Its difference from the piecewise constant interpolation tends to zero
in $L^2((a,b);V^*)$, by integrating that derivative on each step.
Testing in time against smooth scalar functions and in space against $V$
identifies the limit derivative as $-B(u^m)$.

For \eqref{eq:scalar}, pair the exact Euler equations directly with $f$.
Self-adjointness gives
$\langle z_j-z_{j-1},f\rangle=-h\langle z_j^m,Bf\rangle$.
The power map is Lipschitz on $[0,K_v]$, so the right sums converge in
$L^1$ against the bounded function $Bf$. An incomplete last time step
has absolute value at most
$h\norm{Bf}_\infty K_v^{m-1}\norm v_1$.
The left side converges against bounded $f$. This proves
\eqref{eq:scalar} at every time, by $L^1$ continuity.
\end{proof}

For later use, if $f\in D(L)$, spectral calculus gives
\begin{equation}\label{eq:fractional}
 Bf=c_s\int_0^\infty(f-P_rf)\frac{dr}{r^{1+s}},
 \qquad c_s=\frac{s}{\Gamma(1-s)}.
\end{equation}
The integral converges in $L^2$. If also $f,Lf$ are bounded, then
\[
 \norm{f-P_rf}_\infty
 \le\min\{2\norm f_\infty,r\norm{Lf}_\infty\}
\]
shows that $Bf$ is bounded. In particular every $C_c^\infty(M)$ function
is permitted in \eqref{eq:scalar}. This also gives the distributional
equation \eqref{eq:weak} for the bounded-data solutions.

Choose nonnegative $v_n\in C_c^\infty(M)$ of integral $A$, supported
in a neighborhood of $o$ of radius tending to zero, all within one
compact set $K_0$. A normalized bump in a fixed coordinate chart gives
such a sequence. Let
\[
 u_n(t)=S_tv_n,\qquad w_n=u_n^m,
 \qquad M_A(t)=C_0\max\{A^\beta t^{-\alpha},A\}.
\]
Then $\norm{u_n(t)}_1\le A$ and $\norm{u_n(t)}_\infty\le M_A(t)$.
Since $1-\alpha(m-1)=\beta>0$, we have the uniform estimate
\begin{equation}\label{eq:flux}
 \int_0^t\norm{w_n(r)}_1\,dr\le F_A(t):=
 C_0^{m-1}\left(
 \frac{A^{1+\beta(m-1)}}{\beta}t^\beta+A^mt\right),
 \qquad F_A(t)\longrightarrow0\quad(t\downarrow0).
\end{equation}
For every fixed $0<a<b$, Lemma~\ref{lem:energy} and smoothing also give
\begin{equation}\label{eq:uniformenergy}
 \sup_n\left(\norm{w_n}_{L^2(a,b;V)}+
       \norm{\partial_tu_n}_{L^2(a,b;V^*)}\right)<\infty.
\end{equation}

\section{Spatial tightness from the linear heat flow}

Local convergence on $M$ leaves open the possibility that part of the
approximating mass escapes to infinity. We obtain tightness by testing
with functions whose fractional Laplacian is uniformly small.

\begin{lemma}\label{lem:averages}
Let $0\le\chi_j\le1$ be smooth compactly supported functions increasing
to one. For $H>0$ define
\begin{equation}\label{eq:average}
 f_{H,j}=\frac1H\int_0^H P_r\chi_j\,dr.
\end{equation}
Then $f_{H,j}\in D(L)\cap C_0(M)$, $0\le f_{H,j}\le1$, and
\begin{equation}\label{eq:averagebound}
 \norm{Bf_{H,j}}_\infty\le C_s H^{-s},
 \qquad C_s=c_s\left(\frac1{1-s}+\frac1s\right).
\end{equation}
For fixed $H$, the functions $f_{H,j}$ increase to one uniformly on
each compact subset of $M$.
\end{lemma}

\begin{proof}
In the spectral calculus, the average in \eqref{eq:average} has multiplier
$(1-e^{-H\lambda})/(H\lambda)$, with value one at zero. This multiplier
and its product with $\lambda$ are bounded. Thus
\[
 f_{H,j}\in D(L),\qquad
 Lf_{H,j}=\frac{\chi_j-P_H\chi_j}{H},\qquad
 \norm{Lf_{H,j}}_\infty\le H^{-1}.
\]
Positivity and contraction give $0\le f_{H,j},P_rf_{H,j}\le1$.
The semigroup law gives
\[
 P_rf_{H,j}-f_{H,j}
 =\frac1H\left(\int_H^{H+r}P_\tau\chi_j\,d\tau
                     -\int_0^rP_\tau\chi_j\,d\tau\right),
\]
so $\norm{f_{H,j}-P_rf_{H,j}}_\infty\le\min\{1,r/H\}$.
Splitting \eqref{eq:fractional} at $r=H$ proves
\eqref{eq:averagebound}. In particular these functions belong to the
test class of Lemma~\ref{lem:energy}.

Continuity follows by integrating the continuous heat flow over
$[\eta,H]$ and then letting $\eta\downarrow0$, with uniform error at
most $\eta/H$. If $R=d(x,\supp\chi_j)$, \eqref{eq:heatkernel} gives
\[
 f_{H,j}(x)\le \frac{C\norm{\chi_j}_1}{H}
           \int_0^H r^{-N/2}e^{-R^2/(Cr)}\,dr.
\]
For $R\ge1$, the last integral is at most
\[
 e^{-R^2/(2CH)}\int_0^H r^{-N/2}e^{-1/(2Cr)}\,dr,
\]
which tends to zero as $R\to\infty$. Completeness makes bounded closed
neighborhoods of a compact set compact, so $f_{H,j}\in C_0(M)$.
Finally $P_r1=1$ and monotone convergence give $f_{H,j}\uparrow1$.
For any $\eta>0$, the increasing open sets
$\{f_{H,j}>1-\eta\}$ cover a given compact set; a finite subcover gives
the stated uniform convergence.
\end{proof}

\begin{proposition}\label{prop:tightness}
For every $T<\infty$ and $\varepsilon>0$, there is a compact set
$K\subset M$ such that
\begin{equation}\label{eq:tightness}
 \sup_n\sup_{0\le t\le T}\int_{M\setminus K}u_n(t)\dd\le\varepsilon.
\end{equation}
\end{proposition}

\begin{proof}
Choose $H$ so large that $C_sH^{-s}F_A(T)\le\varepsilon/4$.
Then choose $j$ so large that
$A\sup_{K_0}(1-f_{H,j})\le\varepsilon/4$.
Using \eqref{eq:scalar}, \eqref{eq:flux} and
\eqref{eq:averagebound}, for all $n$ and $t\le T$ we get
\[
 \int_Mu_n(t)(1-f_{H,j})\dd
 \le A-\int_Mu_n(t)f_{H,j}\dd\le\varepsilon/2.
\]
There is a compact $K$ outside which $f_{H,j}\le1/2$.
This proves \eqref{eq:tightness}.
\end{proof}

\section{Compactness and the initial trace}

\begin{lemma}\label{lem:compactness}
A subsequence of $u_n$ converges locally in space-time $L^{2m}$ on
$(0,\infty)\times M$. There is one set of full measure in positive time
on which the same subsequence converges in $L^1$ on every compact subset
of $M$.
\end{lemma}

\begin{proof}
Fix $0<a<b$. From \eqref{eq:uniformenergy}, for $0<h<b-a$,
\[
 \norm{u_n(\cdot+h)-u_n}_{L^2(a,b-h;V^*)}\le Ch.
\]
Pair this difference with $w_n(t+h)-w_n(t)$ in $V$ and use
Cauchy--Schwarz to get
\begin{equation}\label{eq:pairedshift}
 0\le \int_a^{b-h}\!\int_M
 (u_n(t+h)-u_n(t))(w_n(t+h)-w_n(t))\dd\,dt\le Ch.
\end{equation}
The pairing is the usual integral, since $u_n(t)\in L^2(M)$.
On the common range $[0,M_A(a)]$,
\[
 |p^m-q^m|^2\le mM_A(a)^{m-1}(p-q)(p^m-q^m).
\]
Consequently the powers satisfy the global time-translation bound
\begin{equation}\label{eq:timeshift}
 \int_a^{b-h}\norm{w_n(t+h)-w_n(t)}_2^2\,dt\le Ch.
\end{equation}

For a compact set $K$ and $\delta>0$, the heat map
$P_\delta:L^2(M)\to L^2(K)$ has square-integrable kernel, because
\[
 \int_K\int_Mp_\delta(x,y)^2\dd(y)\dd(x)
       =\int_Kp_{2\delta}(x,x)\dd(x)<\infty.
\]
It can therefore be approximated in operator norm by finite-rank maps:
approximate the kernel in product $L^2$ by finite sums of separated
functions and apply Cauchy--Schwarz. On the other hand, the scalar
inequality $|1-e^{-x}|^2\le C_sx^s$ gives
\begin{equation}\label{eq:heatsmooth}
 \norm{v-P_\delta v}_2^2\le C_s\delta^s\E(v,v),\qquad v\in V.
\end{equation}
Thus $w_n$ is uniformly approximated in $L^2((a,b)\times K)$ by the
heat-smoothed family and then by a finite-rank family. Its coefficients
inherit \eqref{eq:timeshift}, since they are bounded linear functionals
on the global space $L^2(M)$.

For completeness, a coefficient vector $q(t)$ can be replaced by its
averages $q_I$ on equal time intervals of length $\ell$ dividing $b-a$.
Its squared approximation error satisfies
\[
 \sum_I\int_I|q-q_I|^2
 =\sum_I\frac1{2\ell}\int_I\int_I|q(t)-q(r)|^2\,dt\,dr
 \le\frac1\ell\int_0^\ell\int_a^{b-h}|q(t+h)-q(t)|^2\,dt\,dh
 \le C\ell.
\]
For a fixed partition these averages form a bounded finite-dimensional
family. Choosing successively $\delta$, spatial rank and $\ell$ proves
local space-time $L^2$ precompactness of $w_n$.

A diagonal argument over compact spatial sets and positive compact time
intervals gives $w_n\to w$ locally in $L^2$ and almost everywhere.
The inequality $|p-q|^m\le|p^m-q^m|$ for $p,q\ge0$ shows that
$u_n\to u=w^{1/m}$ locally in $L^{2m}$.
Choose a further subsequence with summable $L^1$ errors on every member
of a countable space-time exhaustion. Fubini then gives a single
full-measure set of times with local spatial $L^1$ convergence on all
compact sets. Intersecting with the full-measure set of time slices of
the almost-everywhere space-time convergence preserves this property.
\end{proof}

\begin{proof}[Proof of Theorem~\ref{thm:main}]
At every good time $t$ of Lemma~\ref{lem:compactness}, local $L^1$
convergence and Proposition~\ref{prop:tightness} imply global $L^1$
convergence. Indeed, Fatou gives the same small tail for the limit, and
the two tails can be made arbitrarily small before using convergence on
the remaining compact set.
Choose good times $a_j\downarrow0$. The semigroup property and
\eqref{eq:contraction} give
\[
 \sup_{t\ge a_j}\norm{u_n(t)-u_\ell(t)}_1
 \le\norm{u_n(a_j)-u_\ell(a_j)}_1\longrightarrow0.
\]
Completeness of $C([a_j,T];L^1(M))$ therefore gives consistent limits
on overlapping intervals. They define one curve
$u\in C((0,\infty);L^1(M))$, agreeing with the local limit, and
\begin{equation}\label{eq:globalconv}
 u_n\longrightarrow u\quad\text{uniformly in }L^1(M)
            \text{ on every }[a,T]\subset(0,\infty).
\end{equation}
Nonnegativity, $\norm{u(t)}_1\le A$ and
$\norm{u(t)}_\infty\le M_A(t)$ pass to this representative for every
$t>0$. On $[a,T]$ the power map has a common Lipschitz constant, so
\eqref{eq:globalconv} also holds globally in $L^1$ for $u_n^m\to u^m$.

For $\phi\in C_c^\infty((0,\infty)\times M)$, its time support lies
in some $[a,T]$. Equation~\eqref{eq:fractional} gives
\[
 \sup_t\norm{B\phi(t)}_\infty<\infty.
\]
We can therefore pass to the limit
in the approximate distributional equations using the global power
convergence. This proves \eqref{eq:weak}. The right side is absolutely
integrable, since $\norm{u(t)^m}_1\le A M_A(a)^{m-1}$ on $[a,T]$.
The same bounds give all the local integrability and boundedness in
\eqref{eq:regularity}.

For a fixed $\psi\in C_c^\infty(M)$, \eqref{eq:scalar} and
\eqref{eq:flux} imply
\[
 \left|\int_Mu_n(t)\psi\dd-\int_Mv_n\psi\dd\right|
       \le\norm{B\psi}_\infty F_A(t).
\]
Let $n\to\infty$ at fixed $t>0$, using \eqref{eq:globalconv} and the
initial concentration at $o$. Then let $t\downarrow0$. Since
$F_A(t)\to0$, we obtain
$\int_Mu(t)\psi\dd\to A\psi(o)$.
Uniform approximation of compact continuous functions by smooth compact
functions extends this trace to $C_c(M)$, because the mass is bounded by
$A$ at every positive time.

Fix $0\le\chi\le1$, $\chi\in C_c^\infty(M)$, equal to one near $o$.
The compact-test trace and positivity give
\[
 \int_Mu(t)\chi\dd\longrightarrow A,\qquad
 0\le\int_Mu(t)(1-\chi)\dd
       \le A-\int_Mu(t)\chi\dd\longrightarrow0.
\]
For any prescribed $\psi\in C_b(M)$, apply the compact-continuous trace
to $\psi\chi$ and bound the remaining integral by
$\norm\psi_\infty\int_Mu(t)(1-\chi)\dd$. This proves
\eqref{eq:trace} as an actual limit. The same inequalities give
$\norm{u(t)}_1\to A$ as $t\downarrow0$; together with its upper bound
this proves \eqref{eq:mass} for every $T>0$.
\end{proof}

\paragraph{AI use.}
OpenAI Codex assisted with proof development, source checks, manuscript
drafting and internal AI review. The mathematical claims are submitted for
independent review.

\begin{thebibliography}{9}
\bibitem{BBGM}
E.~Berchio, M.~Bonforte, G.~Grillo and M.~Muratori,
\emph{The fractional porous medium equation on noncompact Riemannian manifolds},
Mathematische Annalen \textbf{389} (2024), 3603--3651.
\href{https://doi.org/10.1007/s00208-023-02731-6}{doi:10.1007/s00208-023-02731-6}.
The numbered source statements used here are in
\href{https://arxiv.org/abs/2109.10732v2}{arXiv:2109.10732v2} (2023).

\bibitem{Gri}
A.~Grigor'yan,
\emph{Analytic and geometric background of recurrence and non-explosion
of the Brownian motion on Riemannian manifolds},
Bulletin of the American Mathematical Society \textbf{36} (1999), 135--249.
\href{https://www.math.uni-bielefeld.de/~grigor/recor.pdf}{Author's version}.

\bibitem{GMP}
G.~Grillo, M.~Muratori and F.~Punzo,
\emph{Fractional porous media equations: existence and uniqueness of weak
solutions with measure data},
Calculus of Variations and Partial Differential Equations \textbf{54}
(2015), 3303--3335.
\href{https://doi.org/10.1007/s00526-015-0904-4}{doi:10.1007/s00526-015-0904-4}.
\href{https://arxiv.org/abs/1312.6076v3}{arXiv:1312.6076v3}.

\bibitem{OWR}
G.~Grillo,
\emph{Fractional nonlinear diffusions on manifolds: well-posedness and
smoothing effects}, contribution to \emph{Mini-Workshop: Hardy Inequalities
in Discrete and Continuum Settings}, Oberwolfach Report 14/2025,
597--598.
\href{https://ems.press/content/serial-article-files/51358?nt=1}{Workshop report}.
\end{thebibliography}
\end{document}
